% ============================================================================
\chapter{Uklapanje cijele igre sa video izlazom}
\label{ch:uklapanje}
% ============================================================================

U ovom poglavlju obra\dj{}uje se sredi\v{s}nji in\v{z}enjerski problem rada, a to
je uklapanje cijele igre zajedno sa operativnim sistemom i video izlazom u
stroge resurse plo\v{c}e Tang Nano 9K. Najprije se opisuje problem koji se javio
pri prvom poku\v{s}aju, zatim se analizira njegov stvarni uzrok kroz zauze\'ce
blokovske memorije. Nakon toga prikazuju se dva rje\v{s}enja koja su zajedno
omogu\'cila uklapanje, kao i jedno va\v{z}no otkri\'ce o pona\v{s}anju alata za
sintezu. Na kraju se daje kona\v{c}ni rezultat uklapanja i prilagodba video
inicijalizacije. Ovo poglavlje predstavlja te\v{z}i\v{s}te tehni\v{c}kog
doprinosa rada.

\section{Problem uklapanja}
\label{sec:uklapanje-problem}

Pri prvom poku\v{s}aju spajanja procesora sa operativnim sistemom i video
jedinicom u rezoluciji visoke definicije, sinteza nije uspjela. Procijenjeni
broj logi\v{c}kih elemenata iznosio je oko \v{s}esnaest hiljada, dok plo\v{c}a
raspola\v{z}e sa oko osam i po hiljada. Na prvi pogled \v{c}inilo se da je
problem u koli\v{c}ini logike, me\dj{}utim, dublja analiza pokazala je da je
stvarni uzrok bio drugdje.

\section{Analiza zauze\'ca blokovske memorije}
\label{sec:uklapanje-analiza}

Pregledom izvje\v{s}taja sinteze utvr\dj{}eno je da je ve\'c sama ina\v{c}ica sa
serijskim izlazom tro\v{s}ila svu raspolo\v{z}ivu blokovsku memoriju, odnosno
svih dvadeset \v{s}est blokova. Od toga je programska memorija zauzimala
\v{s}esnaest blokova, radna memorija osam, a serijska jedinica dva. Kako vi\v{s}e
nije bilo slobodne blokovske memorije, video jedinica je svoju znakovnu i mapnu
memoriju morala smjestiti u logi\v{c}ke elemente i registre, \v{s}to je dovelo do
prividne eksplozije broja logi\v{c}kih elemenata. Drugim rije\v{c}ima, problem
nije bio u logici, nego u nedostatku blokovske memorije.

Pri poku\v{s}aju smanjenja memorija uo\v{c}eno je i va\v{z}no svojstvo alata za
sintezu. Alat pouzdano prepoznaje i smje\v{s}ta memoriju u blokovsku memoriju
samo kada je njena dubina jednaka stepenu broja dva. Kada je programska memorija
poku\v{s}ana smanjiti na dubinu koja nije stepen broja dva, alat je memoriju
po\v{c}eo realizovati pomo\'cu registara, \v{s}to je dovelo do gotovo devetnaest
hiljada registara i nove neuspjele sinteze. Ovo otkri\'ce odredilo je smjer
svih daljih optimizacija, jer su sve dubine memorija morale ostati stepeni broja
dva.

\section{Dva rje\v{s}enja: manji program i manja video memorija}
\label{sec:uklapanje-rjesenja}

Prvo rje\v{s}enje odnosi se na programsku memoriju. Utvr\dj{}eno je da je program
povezivao i one funkcije operativnog sistema koje se uop\v{s}te ne koriste. Uz
prevo\dj{}enje sa razdvajanjem funkcija u zasebne sekcije i uklanjanjem
nekori\v{s}tenih sekcija pri povezivanju, veli\v{c}ina programa smanjena je sa
oko dvadeset dva i po hiljade bajta na oko petnaest hiljada bajta. Time je
program pao ispod granice od \v{s}esnaest kilobajta, pa programska memorija sada
zauzima osam umjesto \v{s}esnaest blokova.

Drugo rje\v{s}enje odnosi se na video memoriju. Umjesto rezolucije visoke
definicije, kori\v{s}tena je ni\v{z}a rezolucija, \v{c}ime se znakovna memorija
smanjila sa dvanaest hiljada dvije stotine osamdeset osam na \v{c}etiri hiljade
devedeset \v{s}est rije\v{c}i. Obje veli\v{c}ine ostale su stepeni broja dva,
\v{s}to je u skladu sa ranije opisanim svojstvom alata. Zajedno, ova dva
rje\v{s}enja oslobodila su dovoljno blokovske memorije da se cijeli sistem
uklopi.

\section{Rezultat uklapanja i prilagodba video inicijalizacije}
\label{sec:uklapanje-rezultat}

Nakon primjene oba rje\v{s}enja, sinteza je uspjela, a sistem se uklopio u
resurse plo\v{c}e sa odre\dj{}enom rezervom. Pore\dj{}enje zauze\'ca resursa prije
i poslije optimizacije dato je u tabeli \ref{table:resursi}.

\begin{table}[H]
\caption{Zauze\'ce resursa kombinovanog sistema prije i poslije optimizacije.}
\label{table:resursi}
\begin{center}
\begin{tabular}{lcc}
\hline
Resurs & Prije & Poslije \\ \hline
Logi\v{c}ki elementi & 16041 (prekora\v{c}eno) & 3348 / 8640 \\
Registri & -- & 1349 / 6693 \\
Blokovska memorija & 26 / 26 & 22 / 26 \\
\quad programska (pROM) & 16 & 8 \\
\quad radna (SP) & 8 & 8 \\
\quad ostalo (SDPB) & 2 & 6 \\ \hline
\end{tabular}
\end{center}
\end{table}

Nakon \v{s}to je sklop uklopljen, prilago\dj{}ena je i video inicijalizacija u
programu. Po\v{s}to je mapa polja ista u obje rezolucije, jer se polje sastoji od
istog broja velikih znakova, mijenjala se samo veli\v{c}ina pojedinog znaka.
Inicijalizacija znakova prilago\dj{}ena je da upisuje znakove veli\v{c}ine
\v{s}esnaest sa \v{s}esnaest ta\v{c}aka umjesto ranijih dvadeset \v{c}etiri sa
trideset dvije. Sistem se nakon toga uspje\v{s}no pokre\'ce na hardveru i video
inicijalizacija se izvr\v{s}ava do kraja, dok kona\v{c}na vizuelna potvrda slike
na ekranu ostaje stavka za daljnji rad.

\section{Zaklju\v{c}ak poglavlja}
\label{sec:uklapanje-zakljucak}

U poglavlju je rije\v{s}en glavni in\v{z}enjerski problem rada. Pokazano je da
uzrok neuspjeha prvog uklapanja nije bila koli\v{c}ina logike, nego nedostatak
blokovske memorije, te da alat za sintezu pouzdano koristi blokovsku memoriju
samo za dubine koje su stepeni broja dva. Smanjenjem programa uklanjanjem
nekori\v{s}tenih funkcija i kori\v{s}tenjem ni\v{z}e video rezolucije sistem je
uklopljen u resurse plo\v{c}e i pokrenut na hardveru. Time je pokazano da je
cijela igra sa operativnim sistemom i video izlazom izvodljiva na ovoj
niskobud\v{z}etnoj plo\v{c}i.
